\subsection{Market evidence, plus factors and screens}\label{sec:13-market-evidence}

Detection is the central challenge in the control of coordinated price elevation \citep[p.~256]{Kaplow2013}. Firms
hide conduct that may be illegal, and an aggressive search for it produces false positives that chill competitive
behavior. This subsection asks what an outsider can infer from market evidence before we turn to the specialized
tests for procurement. It draws on two book-length treatments. \citet[chs.~10--11]{Kaplow2013} classifies the
evidence on which a finding of coordination can rest. \citet[chs.~11--12]{MarshallMarx2012} catalog the
plus factors of cartels and bidding rings. A plus factor is economic circumstantial evidence of collusion
beyond the parallel movement of prices (p.~213). \Cref{sec:12-anatomy} introduces the concept and its legal
background. Here we ask how strong each kind of evidence is and how the kinds combine.

\casehead{Parallel prices are not evidence}
Parallel movement of prices is the weakest evidence of all, because competitive markets produce it too. When the
costs of competing firms rise, their prices rise together and their quantities fall. When demand rises, prices and
quantities rise together. Competitors react alike to new technology, new regulations and shifts in tastes
\citep[p.~260]{Kaplow2013}. Kaplow compares parallel pricing to breathing. Every patient breathes, so no doctor
treats breathing as a symptom of bronchitis. Useful evidence is behavior that is consistent with coordination and
inconsistent with competition, or the reverse (p.~261). Kaplow groups such evidence into pricing patterns, measures
of price elevation and facilitating practices. We take these in turn and then consider structural conditions,
internal records and communications.

\casehead{Three phases of a collusive episode}
Kaplow separates three phases: the initiation of elevated prices, their maintenance and their breakdown (p.~261).
Each phase leaves a different pattern. At initiation, a sharp price increase is suspicious only if no concurrent
change in cost or demand can explain it. Any proposed explanation must also match the movement of quantities, as
the next proposition shows. It uses a simple competitive benchmark with price-taking firms.

\begin{proposition}[Sign restrictions on a price increase]\label{prop:13-kaplow-signs}
Let industry demand be $Q = D(p, y)$, with $D_p < 0$ and a demand shifter $y$ with $D_y > 0$ ($D$ and $y$ are
local symbols). Firm $i = 1, \dots, N$, with $N \geq 2$, has marginal cost $c_i(q_i) + \kappa_i$, where $c_i' > 0$
and $\kappa_i$ is a cost shifter. Under competition each firm takes the price as given and supplies
$q_i = s_i(p - \kappa_i)$, where $s_i$ is the inverse of $c_i$ and $s_i' = \frac{1}{c_i'} > 0$. Let
$S' = \sum_i s_i'$. Then, for small changes, the following hold.
\begin{enumerate}[label=(\roman*)]
\item A demand increase raises the price, total output and the output of every firm.
\item A common cost increase $d\kappa$ raises the price by less than $d\kappa$ and lowers total output and the
output of every firm.
\item A cost increase at firm $j$ alone raises the price and lowers total output and the output of firm $j$. It
raises the output of every other firm.
\item A coordinated price increase with $y$ and every $\kappa_i$ unchanged lowers total output. The division of
output among firms is whatever the agreement sets.
\end{enumerate}
\end{proposition}

\begin{proof}
Market clearing requires
\begin{equation*}
  D(p, y) = \sum_{i=1}^{N} s_i(p - \kappa_i) .
\end{equation*}
Differentiate both sides totally:
\begin{equation*}
  D_p\,dp + D_y\,dy = \sum_{i} s_i'\,(dp - d\kappa_i) .
\end{equation*}
Collect the terms in $dp$ and divide by $S' - D_p > 0$:
\begin{equation}\label{eq:13-kaplow-dp}
  dp = \frac{D_y\,dy + \sum_i s_i'\,d\kappa_i}{S' - D_p} .
\end{equation}
(i) With $d\kappa_i = 0$ for all $i$, \eqref{eq:13-kaplow-dp} gives $dp = \frac{D_y\,dy}{S' - D_p} > 0$, and then
$dq_i = s_i'\,dp > 0$ for every $i$.
(ii) With $d\kappa_i = d\kappa > 0$ for all $i$, $dp = \frac{S'}{S' - D_p}\,d\kappa$. The factor lies strictly
between 0 and 1 because $-D_p > 0$, so $0 < dp < d\kappa$. Then $dq_i = s_i'(dp - d\kappa) < 0$ for every $i$.
(iii) With only $d\kappa_j > 0$, $dp = \frac{s_j'}{S' - D_p}\,d\kappa_j$. The factor lies strictly between 0 and 1
because $s_j' < S' - D_p$. Hence $dq_j = s_j'(dp - d\kappa_j) < 0$ and $dq_i = s_i'\,dp > 0$ for $i \neq j$. In
(ii) and (iii) total output changes by $dQ = D_p\,dp < 0$.
(iv) Demand slopes down, so $D(p^1, y) < D(p^0, y)$ for any $p^1 > p^0$.
\end{proof}

\Cref{tab:13-kaplow-signs} collects the predictions. The limit of constant marginal cost, $s_i' \to \infty$,
sharpens them. A demand increase then leaves the price unchanged, and a common cost increase passes through in
full. These are the predictions behind the conduct tests of \cref{ch:conduct} \citep[pp.~273--274]{Kaplow2013}. A
common cost shock is the hardest alternative to rule out, because it moves price and output in the same directions
as coordination. Two checks remain. Cost increases can be measured in input markets, and cost shocks that hit firms
unequally reallocate output under competition (p.~262). Under coordination, firm-specific cost shocks reallocate
less output, because the agreement fixes shares (p.~274).

\begin{table}[htbp]
\centering
\caption{What a price increase implies for quantities under competition and under coordination}
\label{tab:13-kaplow-signs}
\footnotesize
\begin{tabularx}{\textwidth}{@{}>{\raggedright\arraybackslash}p{2.6cm}>{\raggedright\arraybackslash}X>{\raggedright\arraybackslash}p{1.8cm}>{\raggedright\arraybackslash}X>{\raggedright\arraybackslash}X@{}}
\toprule
Source of the increase & Price & Industry output & Output across firms & Check that separates it from coordination\\
\midrule
Demand increase & Rises (unchanged if marginal cost is constant) & Rises & Every firm expands & Output
rises\\
\addlinespace
Common cost increase & Rises by less than the cost (by the full cost if marginal cost is constant) & Falls &
Every firm contracts & Measured input prices\\
\addlinespace
Cost increase at one firm & Rises a little & Falls & Output shifts away from the affected firm & Reallocation of
output\\
\addlinespace
Coordinated elevation & Rises & Falls & Set by the agreement, little reallocation & No matching change in cost or
demand\\
\bottomrule
\end{tabularx}

\smallskip
\begin{minipage}{0.95\textwidth}\footnotesize
Notes: Our summary of \cref{prop:13-kaplow-signs} and of \citet[pp.~260--262 and 273--274]{Kaplow2013}.
\end{minipage}
\end{table}

\casehead{Initiation}
Coordinated increases need not be sudden. Firms that do not fully trust one another may raise prices in small
steps, and a leader may move cautiously while it waits for rivals to follow. Kaplow argues that such signaling can
make coordination easier to infer (pp.~262--263). A sequence of announced moves that cost and demand cannot explain
is more telling than a single jump. Sophisticated firms can also spread increases over time, or time them to
coincide with cost events that are hard to measure (pp.~263--264). The price announcements of the vitamins cartel
show both sides. Compared with the period before 1985, the cartel's announcements were more frequent, came at
regular intervals and involved modest individual increases \citep[pp.~221--222]{MarshallMarx2012}. They were
usually joint, with one firm leading and the others matching within about a week. Firms other than the market
leader, Roche, often led joint announcements after 1985 and rarely did before. About 50 percent of the cartel's
announcements came well before the effective date of the increase, against 5 percent in the earlier benchmark
period (p.~222). Long lead times let the cartel test buyers' acceptance and withdraw a resisted increase before
shares were disturbed.

\casehead{Maintenance}
Coordinated prices, and perhaps market shares, tend to be stickier than competitive ones \citep[pp.~264--265]{Kaplow2013}.
Frequent adjustment requires either negotiation, which raises the risk of detection, or signaling, which risks
being misread. A firm that cuts its price because it sees demand softening may be taken for a cheater and start a
war. Competitors, by contrast, respond to changes promptly. Menu costs, the costs of changing prices as such, also
produce stickiness. They differ from coordination frictions in two ways. Menu costs apply symmetrically to
increases and decreases, and they lead firms to change prices at different times. Sluggishness caused by
coordination has neither feature (p.~265). The bromine record fits this account. Prices changed on average every
41.1 weeks during cooperation and every 11.5 weeks otherwise (\cref{sec:13-bromine}). The variance screen of
\citet{AbrantesMetzEtAl2006}, which flags periods in which prices vary unusually little, rests on the same logic.
\citet{Pesendorfer2000} finds more stable market shares in a bidding ring that did not use sidepayments than in one
that did \citep[p.~264, note~17]{Kaplow2013}.

\casehead{Breakdown}
Sudden, sharp price cuts are as suspicious as sudden increases. If prices fall steeply without an exogenous change
of matching size, the earlier price was a coordinated one \citep[p.~265]{Kaplow2013}. Price wars thus reveal past
collusion, as the railroad and bromine records showed (\cref{sec:13-jec,sec:13-bromine}). Using wars as evidence
creates a danger that the targeting of increases does not. A penalty for price cuts may deter the cuts themselves
and so raise prices. Kaplow replies that the threat of punishment is a necessary condition for collusion, so a
policy that makes wars costly can deter collusion itself (p.~266). Kaplow adds that the problem is more
complicated still (p.~266). Firms might respond
with less sensitive triggers or shorter wars, which leave prices elevated and give buyers rarer relief. Kaplow
therefore suggests leniency, immunity or even rewards for the first price cutter, with full sanctions for the
firms that retaliate (pp.~266--267). The next example shows why one proposal of this kind would fail.

\begin{example}[How large must a price cut be to pay?]\label{ex:13-price-freeze}
\citet[pp.~267--268, note~22]{Kaplow2013} discusses a proposal by Sagi. In selected industries, a price cut
of at least 20 percent would trigger a six-month freeze of all prices, and a cut of at least 10 percent a
three-month freeze. Rivals could not retaliate during the freeze, so the first price cutter would gain. Kaplow's
arithmetic shows how much capacity such a gain requires. Let marginal cost be a constant $c$ and the collusive
price $(1 + m)c$, so that $m$ is the elevation of price above marginal cost. A member sells $q$ at the collusive
price and earns $\pi = mcq$ per period. It cuts its price by the fraction $x$ and multiplies its output by $k$. Its
profit becomes
\begin{equation*}
  \pi^{\mathrm{cut}} = k\,q\,\bigl[(1 - x)(1 + m) - 1\bigr]\,c .
\end{equation*}
Suppose the cut must recover $R$ periods of collusive profit, so that it breaks even when $\pi^{\mathrm{cut}} = R\pi$.
Dividing both sides by $cq$ gives the break-even condition
\begin{equation}\label{eq:13-freeze-breakeven}
  k\bigl[(1 - x)(1 + m) - 1\bigr] = Rm .
\end{equation}
In the short run, $R = 1$. Expand the left side of \eqref{eq:13-freeze-breakeven} and collect the terms in $m$:
\begin{align*}
  k(1 - x)(1 + m) - k &= m ,\\
  m\bigl[k(1 - x) - 1\bigr] &= k - k(1 - x) = kx ,\\
  m &= \frac{kx}{k(1 - x) - 1} .
\end{align*}
With $k = 2$ and $x = 0.2$, $m = \frac{0.4}{1.6 - 1} = \frac{2}{3}$. A firm that can double its output at once breaks
even on a 20 percent cut only if the collusive price exceeds marginal cost by two thirds. The cut takes the price
from $\frac{5}{3}c$ to $\frac{4}{3}c$, which halves the margin while output doubles. With $k = 1.5$, the same formula
gives $m = \frac{0.3}{1.2 - 1} = 1.5$, so the price must exceed marginal cost by 150 percent.

Now suppose a 10 percent cut triggers a three-month freeze, after which collusion collapses and eight quarters of
elevated prices are lost. Kaplow ignores discounting and requires the profit of the freeze quarter to cover the
eight lost quarters, so $R = 8$ and $x = 0.1$. Solving \eqref{eq:13-freeze-breakeven} for $k$ gives
\begin{equation*}
  k = \frac{8m}{0.9(1 + m) - 1} = \frac{8m}{0.9m - 0.1} .
\end{equation*}
The required factor falls as $m$ rises and approaches $\frac{8}{0.9} \approx 8.9$ as $m$ grows large. Kaplow's
``almost ninefold'' output increase is this limit (our reading). The limit treats the margin as falling by the same
10 percent as the price. For finite $m$ the margin falls by more, so the limit is a lower bound. At
$m = \frac{2}{3}$ the requirement is 10.7 (our computation). A firm with a
25 percent share would need a share above 200 percent. Kaplow concludes that the proposal provides nowhere near the
desired incentive to cut prices, and that generous treatment of the first price cutter under ordinary enforcement
is more promising.
\end{example}

\casehead{Evidence of price elevation}
A second family of evidence asks whether current prices exceed marginal cost \citep[pp.~269--275]{Kaplow2013}.
Elevation alone does not show coordination, because a single firm with market power also prices above cost
(p.~269). A direct comparison of price with measured marginal cost is simple in logic and hard in practice. The
analyst must decide which costs are variable and how to allocate joint costs, and average cost is easier to
measure than marginal cost. Capacity creates a trap. Suppose marginal cost is low over most of a firm's output
range and rises steeply near capacity. A firm producing near capacity and pricing at marginal cost then sets a
price well above average cost, and it appears to elevate price when it does not (p.~270). Such errors produce
false positives. Kaplow suggests that enforcers treat more costs as variable and require a larger measured gap.

Comparisons across markets avoid measuring the level of marginal cost. Price is at least marginal cost in every
market, so a higher price in one market is supracompetitive to the extent that cost differences do not explain it
(p.~271). Price discrimination gives a sharper bound, as the next remark shows.

\begin{remark}[Price discrimination bounds the margin]\label{rem:13-discrimination-bound}
Suppose a seller with marginal cost $mc$ sells the same good to two groups of buyers at prices $p_L < p_H$, and
serving the two groups costs the same. The seller would not sell to the low-price group below cost, so
$p_L \geq mc$. Subtract both sides from $p_H$ and divide by $p_H$:
\begin{equation*}
  p_H - mc \;\geq\; p_H - p_L , \qquad \frac{p_H - mc}{p_H} \;\geq\; \frac{p_H - p_L}{p_H} .
\end{equation*}
The Lerner index of the high-price sale is at least the proportional price gap \citep[pp.~271--272]{Kaplow2013}.
The bound applies to negotiated discounts as well as to posted prices. Buyers who obtain large discounts reveal that
other buyers pay well above marginal cost (p.~272, note~33). Two caveats apply. Some price dispersion arises from
buyers' imperfect information, and a highly successful cartel may even reduce dispersion by suppressing secret
discounts (p.~272, note~34).
\end{remark}

Two further sources complete the list. Econometric conduct tests ask whether prices respond to shifts in demand
and cost as competition or as monopoly would predict (pp.~273--274). \Cref{ch:conduct} develops these tests, and
\cref{prop:13-obs-equiv} states their main limit. Entry also carries information (pp.~274--275). If fringe firms
enter or expand while incumbents have unused capacity, prices are probably elevated. A steep fall in price and a
rise in output after entry reveal earlier supracompetitive pricing, much as a price war does. A failure of price to
fall when a key patent expires may also indicate coordination (p.~275, note~45).

\casehead{Facilitating practices}
A facilitating practice (\cref{sec:13-intro}) can matter in two ways \citep[pp.~276--277]{Kaplow2013}. Its use
supports an inference of coordination when no other plausible explanation for it exists. The practice can also be
prohibited itself, which punishes attempts when success is hard to prove. Kaplow examines three groups of practices
(pp.~277--285). The first improves mutual understanding through advance price announcements, exchanges of prices
through trade associations and audits of invoices. More than 150 price-exchange agreements arose in Britain within
a few years after overt price fixing was declared illegal (p.~278, note~53). Kaplow stresses the firms' own
incentives. Helping buyers find discounts would undercut a cartel, while helping sellers spot cheaters would
sustain it. If firms in an industry conducive to coordination make prices more open, the second motive is the
likely one (p.~279). Exchanges of demand forecasts deserve similar scrutiny where transaction prices are secret,
because demand uncertainty is then the main threat to a cartel (p.~281). The second group standardizes products,
for example through basing-point pricing (p.~282). The third changes the payoffs from cutting prices. It includes
cross-ownership, sidepayments, purchases from rivals that sold less than their share, most-favored-customer
clauses, resale price maintenance and meeting-competition clauses (pp.~283--285). \Cref{sec:12-facilitating}
analyzes several of these practices. The publication of concrete prices in \cref{sec:13-concrete}, the
fare-publishing system of \cref{sec:13-atpco} and the interfirm purchases of \cref{sec:13-imperfect-evidence} are
cases in point.

\casehead{Conducive conditions}
Evidence on whether market conditions favor coordination sharpens inference in both directions
\citep[pp.~286--288]{Kaplow2013}. Conduciveness is roughly a necessary condition for success. If conditions make
coordination nearly impossible, it is probably not taking place. The enforcer may misjudge the conditions,
however, and the firms usually understand their industry better. Strong evidence of successful coordination should
therefore override a judgment that conditions are unfavorable (p.~287). Conduciveness is not sufficient either. If
deterrence works, coordination is rare even where conditions are ideal (p.~288). Kaplow reviews the familiar
conditions briefly (pp.~289--292). A larger number of firms makes it harder to agree on a price and a punishment
and harder to detect cheating by small firms. In most prosecuted cases the firms were few and concentration high. Spare capacity raises both the gain from cheating and the ability to punish, while asymmetric capacities
hinder coordination. Many small buyers facing public prices make cheating easy to detect. A few large buyers who
place infrequent, privately negotiated orders make it hard. Public procurement usually publishes winning bids, which
may explain why so many prosecutions involve bid rigging (p.~290). Uncertain demand hides cheating when prices are
secret, and anticipated changes in demand alter the ratio of current gains from cheating to future losses
(\cref{sec:12-demand}). Product heterogeneity complicates agreement and the policing of quality and terms, and most
prosecuted cases involve homogeneous products. Similar costs, horizons and managerial backgrounds ease agreement,
and multimarket contact enlarges the scope for punishment. Finally, coordination pays only if the industry demand
elasticity and the supply response of fringe firms leave room to raise price.

The empirical basis for these factors is weaker than their familiarity suggests (pp.~292--294). Cross-section
studies find little relation between concentration and profitability. Studies within an industry find a more
robust positive relation between concentration and price, although cost differences and endogeneity complicate
them. Structural screens have a poor record. Scherer described a Federal Trade Commission effort to find
price-fixing cases from structural indicators alone as ``a resounding flop'' (pp.~294--295, note~26). Kaplow concludes
that highly unconducive conditions strongly negate success, that strong evidence of success should usually
convince, and that high conduciveness alone is a weak basis for inferring coordination (pp.~294--295). This ranking
differs from that of a rule based on communications. Under such a rule, highly conducive conditions can lower the
inference that prohibited communications were used, because the firms needed them less (p.~286, note~1).
\Cref{sec:12-kaplow} calls this the paradox of proof.

The Sherbrooke cartel of \cref{sec:13-erutku} shows why the checklist deserves limited weight.
\citet[Table~3]{ErutkuHildebrand2010} score the market against eight items derived from \citet{Stigler1964}, and
only three point toward collusion. The market had about 60 stations and a Herfindahl index below 1,800. Yet at
least 80 percent of the stations took part in the cartel. In Kaplow's terms, the wiretaps supplied strong evidence
of success, which outweighed an unfavorable reading of the conditions. The case fits his warning that outsiders
misjudge conditions that the firms themselves understand (p.~287).

\casehead{Internal evidence and communications}
For almost every market-based indicator there is matching evidence inside the firm \citep[pp.~295--302]{Kaplow2013}.
Invoices record prices and quantities. Cost and marketing data reveal what the firm believes about marginal cost
and demand. Strategy memos and notes of meetings explain why prices moved. Omissions are informative too. If a
firm's files contain no new data or discussion about demand, a demand shift is an unlikely explanation for its
price increase (p.~299). Internal evidence has two limits. A large firm's knowledge is spread over many people who
may disagree, and lawyers who anticipate litigation can sanitize records (pp.~299--300). Kaplow doubts that
sanitizing can go far. Several firms would have to fabricate consistent records of a cost shock that never
occurred, quickly and without leaving traces, and their story would have to match external data (p.~301).

Communications among firms are internal evidence of a special kind (pp.~302--306). Explicit and reliable records
of them can be decisive, especially when they refer to earlier successful price increases. Public communications,
such as sequences of advance announcements that converge on a common price, can be as revealing as secret ones
(p.~303). The absence of such records proves little, because firms hide their communications and need them least
where conditions are most conducive (p.~305). Kaplow also argues that inferring secret communications from
circumstantial evidence adds nothing. Let $C$ denote successful coordination, $E$ the observed evidence and $M$ an
unobserved communication (local symbols). By the law of total probability,
\begin{equation*}
  \Pr(C \mid E) = \Pr(C \mid E, M)\,\Pr(M \mid E) + \Pr(C \mid E, \text{not } M)\,\Pr(\text{not } M \mid E) .
\end{equation*}
Any inference about $M$ drawn from $E$ is already part of $\Pr(C \mid E)$. Only direct evidence of $M$, which
enlarges $E$, can strengthen the conclusion (p.~305, note~41).

\casehead{What decided the cases of this chapter}
Kaplow's categories classify the evidence behind the cases of \cref{tab:13-discovery}. The legal cartels, namely the
Joint Executive Committee, the Sugar Institute and the bromine pools, left internal records and open communications.
Wiretapped communications and guilty pleas proved the Quebec gasoline cartel. The lysine and citric acid cartels
were proved in criminal prosecutions and Commission decisions resting on documents and testimony, and the lysine
investigation recorded cartel meetings on video \citep[p.~197, note~46]{AskerNocke2021}. The stamp ring's knockout
records were internal evidence, found after a complaint. The Overlap and fare-publishing cases targeted the
agreements and information systems themselves. The Danish concrete episode is the one case in which market-based
evidence carried the main weight, and there the authority itself had created the transparency. Statistical tests
decided none of the prosecutions. On Long Island and in Ohio they corroborated convictions and confessions after
the fact. Porter and Zona call their tests ``poor substitutes for a wiretap or a disclosure by a dissident ring
member'' \citep[p.~210, note~78]{AskerNocke2021}.

\casehead{Plus factors}
Courts that require proof of an agreement look for plus factors (\cref{sec:12-anatomy}). A plus factor that
delivers a strong inference of collusion is a \emph{super-plus factor} \citep[p.~213]{MarshallMarx2012}. Marshall
and Marx organize plus factors by the cartel conduct that produces them, in nine groups (pp.~213--214). Pricing
structures produce price elevation, quantity restriction, steps to reduce buyers' resistance to price increases and
changes in incentives within firms. Allocation structures produce an allocation of the collusive gain and
redistributions among members. Enforcement structures produce communication and monitoring and a readiness to
abandon collusion if cheating persists. A cartel that has suppressed rivalry may finally engage in dominant-firm
conduct. Marshall and Marx measure the strength of a plus factor $F$ by the likelihood ratio
$S = \frac{\Pr[F \mid C]}{\Pr[F \mid \text{not } C]}$, where $C$ is again the event that the firms collude (p.~235). The
next remark derives the posterior probability of collusion and extends it to several plus factors.

\begin{remark}[Combining plus factors]\label{rem:13-plus-bayes}
Let $O = \frac{\Pr[\text{not } C]}{\Pr[C]}$ be the prior odds against collusion. Bayes' rule gives the posterior
odds in favor of collusion, which we then invert:
\begin{align*}
  \frac{\Pr[C \mid F]}{\Pr[\text{not } C \mid F]}
    &= \frac{\Pr[F \mid C]}{\Pr[F \mid \text{not } C]}\cdot\frac{\Pr[C]}{\Pr[\text{not } C]} = \frac{S}{O} ,\\
  \Pr[C \mid F] &= \frac{\frac{S}{O}}{1 + \frac{S}{O}} = \frac{1}{1 + \frac{O}{S}} .
\end{align*}
This is the formula of \citet[p.~235, note~42]{MarshallMarx2012}, derived as \eqref{eq:12-plus-strength} in
\cref{ch:collusion-theory}. Now suppose that $K$ plus factors $F_1, \dots, F_K$ are independent conditional on $C$
and conditional on not $C$. The likelihood ratio of the joint event is then the product of the individual ratios,
$S = \prod_{k=1}^{K} S_k$. As an illustration with our own numbers, let each plus factor have
$\Pr[F_k \mid C] = 0.6$ and $\Pr[F_k \mid \text{not } C] = 0.2$, so that $S_k = 3$, and let $\Pr[C] = 0.1$, so that
$O = 9$. One plus factor gives $\Pr[C \mid F] = \frac{1}{1 + 3} = 0.25$. Two give $\frac{1}{1 + 1} = 0.5$, and three
give $\frac{1}{1 + \frac{1}{3}} = 0.75$. Plus factors that are weak alone can thus be strong together, as Marshall
and Marx note (p.~235). Conditional independence is a strong assumption, however. Plus factors that share a cause,
such as a concentrated market structure, carry less joint information than the product suggests. Market structure
belongs in the prior odds $O$, which is where Kaplow's conducive conditions enter.
\end{remark}

Marshall and Marx list six super-plus factors (pp.~235--236), which \cref{sec:12-anatomy} reproduces. Two of them
concern market data directly. The first is a path of transaction prices that departs from the predictions of a
reliable econometric model fitted to noncollusive benchmark data. The sixth is stable market shares, customer
incumbency or geographic dominance in the face of excess capacity, while prices and profits rise. The others are
communication and information sharing, interfirm transfers, synchronized changes in sales-force incentives, and
dominant-firm conduct by firms that lack dominance. Cartels respond to detection by avoiding super-plus factors,
and avoidance is costly (pp.~236--237). A market-share cartel that forgoes year-end transfers must police shares more
tightly. One that fakes volatility in its shares can destabilize itself.

\casehead{Plus factors for bidding rings}
\citet[ch.~12]{MarshallMarx2012} adapt the analysis to auctions and procurements. Two cautions come first
(pp.~241--243). An auctioneer or procurement officer who is bribed, or who simply favors one bidder, can produce
depressed prices and skewed awards that look like collusion among bidders. A lone bidder can also enlist shill
bidders to satisfy a rule that requires a minimum number of bids. Next, a ring with $n_{\mathrm{ring}}$ of the $n$
bidders has $n_{\mathrm{ring}} - 1$ free bids to use as it wishes. Screens that look for identical bids therefore presume a careless ring
(p.~243). Two super-plus factors apply to both auction formats (p.~244). One is transfers among bidders, such as
knockout payments, or subcontracting and joint ownership that were not declared before the auction. The other is a
divergence between realized prices and the predictions of a model fitted to periods without a ring.

Other plus factors depend on the format. At an ascending auction, a ring's simplest rule is that members do not bid
against an active member (pp.~244--246). The winner and the last underbidder are then never both members. A group of
bidders who never bid against one another at meaningful prices is suspicious, although many auctions are needed to
see the pattern. Resale of the item soon after the auction to another bidder who was present signals a post-auction
knockout. An ascending auction allocates efficiently, so a ring would have to change its mechanism to avoid this
footprint. Light attendance is no super-plus factor, because some rings require members to attend. At a
sealed-bid auction, post-auction resale is weak evidence, since asymmetric bidders make the auction inefficient
even without a ring (p.~246). Losing bids unrelated to costs are evidence that a ring can avoid. The ring can ask
its members to submit rescaled versions of their noncollusive bids, all below the designated bid (pp.~246--247).
This is the adaptation that \cref{sec:13-pz93} discusses for the rank test.

Close pairs of bids are harder to avoid (p.~247). A designated bidder who bids far below its noncollusive bid tempts
other members to win secretly through a shill. In the example of Marshall and Marx, the highest-valuing member
values the item at \$100,000 and would bid \$90,000 alone, but the ring wants it to bid \$50,000. A member with
value \$60,000 whose share of the collusive gain is below \$10,000 would then gain by bidding about \$50,001 through
a shill. The ring therefore instructs the designated member to bid more, say \$65,000, and has another member bid
just below that amount. The second bid ensures that the designated member complies, because a lower bid would lose
to it. Frequent close pairs of bids then point to collusion, a pattern related to the clustering of
\cref{ex:13-gap-uniform}. Finally, a switch from ascending to sealed bidding that raises prices sharply indicates a
ring at the ascending auctions. Risk aversion also raises sealed-bid prices, but by much less (p.~248).

\begin{example}[Shill bidders in Russian oil and gas auctions]\label{ex:13-russia}
Russian ascending auctions for oil and gas leases required at least two registered bidders
\citep[pp.~248--255]{MarshallMarx2012}. The government set a reserve price and the bid increment, and the lowest
feasible bid was one increment above the reserve. Marshall and Marx have data on 620 auctions from 2004 to 2007,
complete for 121 of them. Five of these were canceled and one used criteria other than price, which leaves 115.
\Cref{tab:13-mm-russia} reproduces their Table~12.2. Auctions with two participants look different from the rest.
Of the 43 such auctions, 34 (79 percent) ended one or two increments above the reserve. With six or more
participants, every winning bid was at least 20 increments above the reserve. The average reserve price of the 34
low-priced two-bidder auctions was 27,515,294 rubles, about half the average over all 620 auctions. The authors
write ``less than half'' (p.~250). Their averages on p.~249 put the ratio at 51 percent (our computation).
Small
licenses of only regional interest could explain part of the pattern. The identities of the losers suggest more
(their Table~12.3). The 34 auctions have 24 distinct winners and 26 distinct losers. Bidder Y won six of them, each
time against bidder 1Y, which never won any of the 620 auctions. Marshall and Marx classify 13 losers as pure
shills, which bid at most once, never won a lease and never took part in an auction without their winning partner.
Two losers rotated as winner and loser, and the remaining 11 are occasional shills. The evidence fits single real
bidders who registered agents to satisfy the two-bidder rule (pp.~251--253).
\end{example}

\begin{table}[htbp]
\centering
\caption{Winning bids in Russian oil and gas lease auctions, by number of participants}\label{tab:13-mm-russia}
\small
\begin{tabular}{@{}lrrrrrr@{}}
\toprule
 & \multicolumn{5}{c}{Increments of the winning bid above the reserve} & \\
\cmidrule(lr){2-6}
Participants & 1 & 2 & 3--19 & 20--99 & 100 or more & Auctions\\
\midrule
2 & 53\% & 26\% & 9\% & 7\% & 5\% & 43\\
3 & 3\% & 9\% & 26\% & 41\% & 21\% & 34\\
4 & 8\% & 0\% & 15\% & 15\% & 62\% & 13\\
5 & 10\% & 10\% & 20\% & 10\% & 50\% & 10\\
6 or more & 0\% & 0\% & 0\% & 33\% & 67\% & 15\\
\addlinespace
All & 23\% & 13\% & 15\% & 22\% & 28\% & 115\\
\bottomrule
\end{tabular}

\smallskip
\begin{minipage}{0.86\textwidth}\footnotesize
Notes: Shares of the 115 auctions with complete data, 2004--2007. Rows sum to 100 percent up to rounding. Source:
\citet{MarshallMarx2012}, Table~12.2.
\end{minipage}
\end{table}

The remainder of this section turns some of these plus factors into formal statistical procedures for
procurement.
